\subsection{DECOMPOSITION OF TENSOR PRODUCTS OF SIMPLE g-MODULES}

PROPOSITION 2. Let \(\lambda, \mu \in \mathrm{P}_{++}\) . In \(\mathcal{R}(\mathfrak{g})\) , we have

\[
[ \lambda ]. [ \mu ] = \sum_ {\nu \in \mathrm{P} _ {+ +}} m (\lambda , \mu , \nu) [ \nu ]
\]

with

\[
m (\lambda , \mu , \nu) = \sum_ {w, w ^ {\prime} \in \mathrm{W}} \varepsilon (w w ^ {\prime}) \mathfrak {P} (w (\lambda + \rho) + w ^ {\prime} (\mu + \rho) - (\nu + 2 \rho)).
\]

Let E, F be finite dimensional simple g-modules of highest weights \(\lambda, \mu\) . Let \(l_{\nu}\) be the length of the isotypical component of \(E \otimes F\) of highest weight \(\nu\) . It suffices to show that

\[
l _ {\nu} = \sum_ {w, w ^ {\prime} \in \mathrm{W}} \varepsilon (w w ^ {\prime}) \mathfrak {P} (w (\lambda + \rho) + w ^ {\prime} (\mu + \rho) - (\nu + 2 \rho)).\tag{2}
\]

Put \(c_{1} = \operatorname{ch}(\mathrm{E}) = \sum_{\sigma \in \mathrm{P}} m_{\sigma} e^{\sigma}\) , \(c_{2} = \operatorname{ch}(\mathrm{F})\) , and \(d = \mathrm{J}(e^{\rho})\) , where \(\mathrm{J}\) is defined as in no. 2. We have

\[
\sum_ {\xi \in \mathrm{P} _ {+ +}} l _ {\xi} \mathrm{ch} [ \xi ] = \mathrm{ch} (\mathrm{E} \otimes \mathrm{F}) = c _ {1} c _ {2}
\]

so, after multiplying by \(d\) and using Th. 1,

\[
\begin{array}{c} \sum_ {\xi \in \mathrm{P} _ {+ +}} l _ {\xi} \mathrm{J} (e ^ {\xi + \rho}) = c _ {1} \mathrm{J} (e ^ {\mu + \rho}) = \left(\sum_ {\sigma \in \mathrm{P}} m _ {\sigma} e ^ {\sigma}\right) \left(\sum_ {w \in \mathrm{W}} \varepsilon (w) e ^ {w (\mu + \rho)}\right) \\ = \sum_ {\tau \in \mathrm{P}} \left(\sum_ {w \in \mathrm{W}} \varepsilon (w) m _ {\tau + \rho - w (\mu + \rho)}\right) e ^ {\tau + \rho}. \end{array}\tag{3}
\]

Now, if \(\xi \in \mathrm{P}_{++}\) , \(\xi + \rho\) belongs to the chamber defined by B (Chap. VI, §1, no. 10); thus, for all \(w \in \mathrm{W}\) distinct from 1, we have \(w(\xi + \rho) \notin \mathrm{P}_{++}\) . Consequently, the coefficient of \(e^{\nu + \rho}\) in \(\sum_{\xi \in \mathrm{P}_{++}} l_{\xi} \mathrm{J}(e^{\xi + \rho})\) is equal to \(l_{\nu}\) . In view of (3), we obtain

\[
l _ {\nu} = \sum_ {w \in \mathrm{W}} \varepsilon (w) m _ {\nu + \rho - w (\mu + \rho)},
\]

that is, by Prop. 1,

\[
l _ {\nu} = \sum_ {w, w ^ {\prime} \in \mathrm{W}} \varepsilon (w) \varepsilon (w ^ {\prime}) \mathfrak {P} (w ^ {\prime} (\lambda + \rho) - (\nu + \rho - w (\mu + \rho) + \rho))
\]

which proves (2).

Example. We return to the Example of no. 1. Let \(\lambda = (n/2)\alpha, \mu = (p/2)\alpha\) , \(\nu = (q/2)\alpha\) with \(n \geq p\) . We have

\[
\begin{array}{l} m (\lambda , \mu , \nu) = \mathfrak {P} \left(\frac {n}{2} \alpha + \frac {\alpha}{2} + \frac {p}{2} \alpha + \frac {\alpha}{2} - \frac {q}{2} \alpha - \alpha\right) \\ \qquad - \mathfrak {P} \left(\frac {n}{2} \alpha + \frac {\alpha}{2} - \frac {p}{2} \alpha - \frac {\alpha}{2} - \frac {q}{2} \alpha - \alpha\right) \\ \qquad - \mathfrak {P} \left(- \frac {n}{2} \alpha - \frac {\alpha}{2} + \frac {p}{2} \alpha + \frac {\alpha}{2} - \frac {q}{2} \alpha - \alpha\right) \\ \qquad + \mathfrak {P} \left(- \frac {n}{2} \alpha - \frac {\alpha}{2} - \frac {p}{2} \alpha - \frac {\alpha}{2} - \frac {q}{2} \alpha - \alpha\right) \\ \qquad = \mathfrak {P} \left(\frac {n + p - q}{2} \alpha\right) - \mathfrak {P} \left(\frac {n - p - q - 2}{2} \alpha\right). \end{array}
\]

This is zero if \(n + p + q\) is not divisible by 2, or if \(q \geq n + p\) . If

\[
q = n + p - 2 r
\]

with \(r\) an integer \(\geq 0\) , we have

\[
m (\lambda , \mu , \nu) = \mathfrak {P} (r \alpha) - \mathfrak {P} ((r - p - 1) \alpha)
\]

hence \(m(\lambda,\mu,\nu)=1\) if \(r\leq p\) and \(m(\lambda,\mu,\nu)=0\) if r\textgreater p.~Finally, the g-module \(\mathrm{V}(n)\otimes\mathrm{V}(p)\) is isomorphic to

\[
\mathrm{V} (n + p) \oplus \mathrm{V} (n + p - 2) \oplus \mathrm{V} (n + p - 4) \oplus \dots \oplus \mathrm{V} (n - p)
\]

(Clebsch-Gordan formula).

